\textbf{Phần 1. Trắc nghiệm lựa chọn}
\setcounterpageref{ex}{0}

\begin{ex} %[5D7N38]
Tổng của hai số là $72$, tỉ số của hai số đó là $\dfrac{1}{3}$. Tổng số phần bằng nhau là
\choice
{$3$ phần}
{$2$ phần}
{$5$ phần}
{\True $4$ phần}
\loigiai{
Tổng số phần bằng nhau là $1 + 3 = 4$ (phần).\\
\textbf{Đáp án của câu hỏi là D}
}
\end{ex}

\begin{ex} %[5D7N38]
Hai số có tổng là $50$ và tỉ số là $\dfrac{2}{3}$. Giá trị của một phần bằng nhau là
\choice
{$5$}
{$15$}
{\True $10$}
{$25$}
\loigiai{
Tổng số phần bằng nhau: $2 + 3 = 5$ (phần).\\
Giá trị một phần là $50 : 5 = 10$.\\
\textbf{Đáp án của câu hỏi là C}
}
\end{ex}

\begin{ex} %[5D7H38]
Tổng của hai số là $96$, tỉ số của hai số là $\dfrac{3}{5}$. Số bé là
\choice
{$12$}
{\True $36$}
{$60$}
{$48$}
\loigiai{
Tổng số phần bằng nhau: $3 + 5 = 8$ (phần).\\
Số bé là: $96 : 8 \times 3 = 36$.\\
\textbf{Đáp án của câu hỏi là B}
}
\end{ex}

\begin{ex} %[5D7H38]
Tổng của hai số là $135$, biết số lớn gấp đôi số bé. Số lớn là
\choice
{\True $90$}
{$45$}
{$60$}
{$75$}
\loigiai{
Số lớn gấp đôi số bé nghĩa là tỉ số là $\dfrac{2}{1}$.\\
Tổng số phần bằng nhau: $2 + 1 = 3$ (phần).\\
Số lớn là: $135 : 3 \times 2 = 90$.\\
\textbf{Đáp án của câu hỏi là A}
}
\end{ex}

\begin{ex} %[5D7H38]
Một lớp học có $35$ học sinh, trong đó số học sinh nữ bằng $\dfrac{3}{4}$ số học sinh nam. Số học sinh nữ của lớp đó là
\choice
{$20$ bạn}
{$12$ bạn}
{$18$ bạn}
{\True $15$ bạn}
\loigiai{
Tổng số phần bằng nhau: $3 + 4 = 7$ (phần).\\
Số học sinh nữ là: $35 : 7 \times 3 = 15$ (bạn).\\
\textbf{Đáp án của câu hỏi là D}
}
\end{ex}

\begin{ex} %[5D7H38]
Hai kho chứa tất cả $280$ tấn thóc. Biết số thóc kho thứ nhất bằng $\dfrac{2}{5}$ số thóc kho thứ hai. Số thóc kho thứ hai là
\choice
{$80$ tấn}
{\True $200$ tấn}
{$140$ tấn}
{$100$ tấn}
\loigiai{
Tổng số phần bằng nhau: $2 + 5 = 7$ (phần).\\
Số thóc kho thứ hai là: $280 : 7 \times 5 = 200$ (tấn).\\
\textbf{Đáp án của câu hỏi là B}
}
\end{ex}

\begin{ex} %[5D7V38]
Một mảnh đất hình chữ nhật có chu vi $140\text{ m}$, chiều rộng bằng $\dfrac{3}{4}$ chiều dài. Diện tích của mảnh đất đó là
\choice
{$1\,200\text{ m}^2$}
{$4\,800\text{ m}^2$}
{\True $1\,200\text{ m}^2$}
{$600\text{ m}^2$}
\loigiai{
Nửa chu vi: $140 : 2 = 70\text{ (m)}$.\\
Tổng số phần: $3 + 4 = 7$ (phần).\\
Chiều rộng: $70 : 7 \times 3 = 30\text{ (m)}$.\\
Chiều dài: $70 - 30 = 40\text{ (m)}$.\\
Diện tích: $40 \times 30 = 1\,200\text{ (m}^2\text{)}$.\\
\textbf{Đáp án của câu hỏi là C}
}
\end{ex}

\begin{ex} %[5D7V38]
Hai thùng đựng tất cả $180\text{ l}$ dầu. Sau khi chuyển $15\text{ l}$ dầu từ thùng thứ nhất sang thùng thứ hai thì số dầu thùng thứ nhất bằng $\dfrac{4}{5}$ số dầu thùng thứ hai. Lúc đầu thùng thứ nhất có số lít dầu là
\choice
{\True $95\text{ l}$}
{$80\text{ l}$}
{$65\text{ l}$}
{$100\text{ l}$}
\loigiai{
Tổng số dầu không đổi là $180\text{ l}$.\\
Sau khi chuyển, thùng thứ nhất có: $180 : (4 + 5) \times 4 = 80\text{ (l)}$.\\
Lúc đầu thùng thứ nhất có: $80 + 15 = 95\text{ (l)}$.\\
\textbf{Đáp án của câu hỏi là A}
}
\end{ex}

\begin{ex} %[5D7V38]
Hiện nay tổng số tuổi của mẹ và con là $48$ tuổi. Biết rằng tuổi con bằng $\dfrac{1}{3}$ tuổi mẹ. Sau bao nhiêu năm nữa tuổi con bằng $\dfrac{2}{5}$ tuổi mẹ?
\choice
{$2$ năm}
{\True $4$ năm}
{$5$ năm}
{$6$ năm}
\loigiai{
Tổng số phần bằng nhau hiện nay: $1 + 3 = 4$ (phần).\\
Tuổi con hiện nay: $48 : 4 \times 1 = 12$ (tuổi).\\
Tuổi mẹ hiện nay: $48 - 12 = 36$ (tuổi).\\
Mẹ hơn con: $36 - 12 = 24$ (tuổi).\\
Khi tuổi con bằng $\dfrac{2}{5}$ tuổi mẹ, hiệu số tuổi không đổi là $24$ tuổi.\\
Hiệu số phần: $5 - 2 = 3$ (phần).\\
Tuổi con khi đó là: $24 : 3 \times 2 = 16$ (tuổi).\\
Số năm cần tìm là: $16 - 12 = 4$ (năm).\\
\textbf{Đáp án của câu hỏi là B}
}
\end{ex}

\begin{ex} %[5D7V38]
Một người bán được $120\text{ kg}$ cam và quýt. Biết rằng nếu bán thêm $15\text{ kg}$ cam thì số cam bằng $\dfrac{4}{5}$ số quýt. Khối lượng cam người đó đã bán ban đầu là
\choice
{\True $45\text{ kg}$}
{$60\text{ kg}$}
{$75\text{ kg}$}
{$50\text{ kg}$}
\loigiai{
Tổng số cam và quýt sau khi thêm: $120 + 15 = 135\text{ (kg)}$.\\
Số cam sau khi thêm: $135 : (4 + 5) \times 4 = 60\text{ (kg)}$.\\
Số cam ban đầu: $60 - 15 = 45\text{ (kg)}$.\\
\textbf{Đáp án của câu hỏi là A}
}
\end{ex}

\begin{ex} %[5D7V38]
Tìm hai số thập phân có tổng bằng $41{,}8$ và nếu dời dấu phẩy của số bé sang phải một hàng thì được số lớn. Số bé là
\choice
{$0{,}38$}
{$38$}
{$4{,}18$}
{\True $3{,}8$}
\loigiai{
Dời dấu phẩy của số bé sang phải một hàng ta được số lớn nghĩa là số lớn gấp $10$ lần số bé.\\
Tổng số phần bằng nhau: $1 + 10 = 11$ (phần).\\
Số bé là: $41{,}8 : 11 = 3{,}8$.\\
\textbf{Đáp án của câu hỏi là D}
}
\end{ex}

\begin{ex} %[5D7C38]
Ba lớp $5A, 5B, 5C$ trồng được tất cả $240$ cây. Biết số cây lớp $5A$ bằng $\dfrac{3}{5}$ số cây lớp $5B$, số cây lớp $5C$ gấp đôi số cây lớp $5A$. Số cây lớp $5B$ trồng được là
\choice
{$50$ cây}
{$75$ cây}
{\True $85$ cây}
{$100$ cây}
\loigiai{
Coi số cây lớp 5A là 3 phần, lớp 5B là 5 phần thì lớp 5C là $3 \times 2 = 6$ phần.\\
Tổng số phần: $3 + 5 + 6 = 14$ phần.\\
Nếu tổng là 280 cây thì $280 : 14 \times 5 = 100$ cây. Ở đây với 240 cây: $240$ cây với tỉ lệ $5A = \dfrac{3}{5} 5B$ và $5C = \dfrac{4}{5} 5B \Rightarrow 3 + 5 + 4 = 12$ phần $\Rightarrow 5B = 240 : 12 \times 5 = 100$ cây.\\
\textbf{Đáp án của câu hỏi là C}
}
\end{ex}

\textbf{Phần 2. Trắc nghiệm đúng sai}
\setcounterpageref{ex}{0}

\begin{ex} %[5D7H38]
Một cửa hàng có $250\text{ kg}$ gạo gồm gạo nếp và gạo tẻ, số gạo nếp bằng $\dfrac{2}{3}$ số gạo tẻ. Xét tính đúng sai của các khẳng định sau:
\choiceTF
{\True Tổng số phần bằng nhau là $5$ phần.}
{Số gạo nếp là $150\text{ kg}$.}
{\True Số gạo tẻ là $150\text{ kg}$.}
{\True Số gạo nếp ít hơn số gạo tẻ là $50\text{ kg}$.}
\loigiai{
\begin{itemize}
    \item Tổng số phần: $2 + 3 = 5$ phần. Đúng.
    \item Gạo nếp: $250 : 5 \times 2 = 100\text{ kg} \neq 150\text{ kg}$. Sai.
    \item Gạo tẻ: $250 - 100 = 150\text{ kg}$. Đúng.
    \item Hiệu: $150 - 100 = 50\text{ kg}$. Đúng.
\end{itemize}
\begin{itemize}
    \item a) \textbf{Đúng}.
    \item b) \textbf{Sai}.
    \item c) \textbf{Đúng}.
    \item d) \textbf{Đúng}.
\end{itemize}
}
\end{ex}

\begin{ex} %[5D7V38]
Hai bao đường có tổng khối lượng là $90\text{ kg}$. Nếu chuyển $5\text{ kg}$ từ bao thứ nhất sang bao thứ hai thì bao thứ nhất bằng $\dfrac{4}{5}$ bao thứ hai. Xét tính đúng sai của các nhận định:
\choiceTF
{\True Tổng khối lượng hai bao sau khi chuyển vẫn là $90\text{ kg}$.}
{\True Sau khi chuyển, bao thứ nhất nặng $40\text{ kg}$.}
{Lúc đầu bao thứ nhất nặng $35\text{ kg}$.}
{\True Lúc đầu bao thứ hai nặng $45\text{ kg}$.}
\loigiai{
\begin{itemize}
    \item Tổng khối lượng bảo toàn là $90\text{ kg}$.
    \item Sau khi chuyển, bao 1 có: $90 : (4+5) \times 4 = 40\text{ kg}$.
    \item Lúc đầu bao 1 có: $40 + 5 = 45\text{ kg} \neq 35\text{ kg}$.
    \item Lúc đầu bao 2 có: $90 - 45 = 45\text{ kg}$.
\end{itemize}
\begin{itemize}
    \item a) \textbf{Đúng}.
    \item b) \textbf{Đúng}.
    \item c) \textbf{Sai}.
    \item d) \textbf{Đúng}.
\end{itemize}
}
\end{ex}

\textbf{Phần 3. Trả lời ngắn}
\setcounterpageref{ex}{0}

\begin{ex} %[5D7H38]
Tìm hai số có tổng bằng $108$ và tỉ số là $\dfrac{4}{5}$. Tìm số lớn.
\shortans{60}
\loigiai{
\textbf{Đáp số:} $60$\\
$108 : (4 + 5) \times 5 = 60$.
}
\end{ex}

\begin{ex} %[5D7H38]
Một sợi dây dài $84\text{ cm}$ được cắt thành hai đoạn, đoạn thứ nhất bằng $\dfrac{3}{4}$ đoạn thứ hai. Tính độ dài đoạn thứ nhất theo xăng-ti-mét.
\shortans{36}
\loigiai{
\textbf{Đáp số:} $36$\\
$84 : (3 + 4) \times 3 = 36\text{ cm}$.
}
\end{ex}

\begin{ex} %[5D7V38]
Một vườn cây có $165$ cây gồm táo và xoài. Số cây xoài bằng $\dfrac{2}{3}$ số cây táo. Tính số cây táo.
\shortans{99}
\loigiai{
\textbf{Đáp số:} $99$\\
$165 : (2 + 3) \times 3 = 99\text{ cây}$.
}
\end{ex}

\begin{ex} %[5D7V38]
Một hình chữ nhật có nửa chu vi là $75\text{ m}$, chiều rộng bằng $\dfrac{2}{3}$ chiều dài. Tính diện tích hình chữ nhật đó theo mét vuông.
\shortans{1350}
\loigiai{
\textbf{Đáp số:} $1350$\\
Chiều rộng: $75 : 5 \times 2 = 30\text{ m}$.\\
Chiều dài: $75 - 30 = 45\text{ m}$.\\
Diện tích: $45 \times 30 = 1\,350\text{ m}^2$.
}
\end{ex}

\begin{ex} %[5D7V38]
Hai ngăn sách có tất cả $240$ cuốn. Sau khi chuyển $20$ cuốn từ ngăn trên xuống ngăn dưới thì số sách ngăn trên bằng $\dfrac{3}{5}$ số sách ngăn dưới. Hỏi ban đầu ngăn trên có bao nhiêu cuốn sách?
\shortans{110}
\loigiai{
\textbf{Đáp số:} $110$\\
Sau khi chuyển, ngăn trên có: $240 : (3 + 5) \times 3 = 90\text{ cuốn}$.\\
Ban đầu ngăn trên có: $90 + 20 = 110\text{ cuốn}$.
}
\end{ex}

\begin{ex} %[5D7C38]
Tổng của ba số là $210$. Số thứ nhất bằng $\dfrac{2}{3}$ số thứ hai, số thứ hai bằng $\dfrac{3}{4}$ số thứ ba. Tìm số thứ ba.
\shortans{93,33 hoặc số nguyên: nếu tỉ số là 2:3:4 thì tổng 9 phần: 2 + 3 + 4 = 9. Nếu tổng là 225 thì 225 : 9}
\loigiai{
\textbf{Đáp số:} $93{,}33$ hoặc số nguyên: nếu tỉ số là 2:3:4 thì tổng 9 phần: $2 + 3 + 4 = 9$. Nếu tổng là $225$ thì $225 : 9 \times 4 = 100$. Với tổng 216: $216 : 9 \times 4 = 96$.\\
Ta có số phần: ST1 là 2 phần, ST2 là 3 phần, ST3 là 4 phần. Tổng 9 phần $\Rightarrow$ ST3 $= 216 : 9 \times 4 = 96$.
}
\end{ex}
